\section{Sujet n\textsuperscript{o}\,26}
\noindent\begin{minipage}{0.5\linewidth}\footnotesize Office du Baccalauréat du Cameroun\end{minipage}%
\hfill\begin{minipage}{0.45\linewidth}\footnotesize\raggedleft Examen : \textbf{Baccalauréat} · Série : \textbf{C/E}\\
Épreuve : \textbf{Mathématiques} · Durée : \textbf{4\,h} · Coef. : \textbf{7}\end{minipage}
\par\smallskip{\color{onbo}\rule{\linewidth}{1pt}}

\noindent\textit{La Camerounaise des Eaux (CDE) modernise le réseau d'eau potable de la ville
de Limbé. Les exercices ci-dessous s'inspirent de cette opération.}

\subsection*{Partie A — Évaluation des ressources \hfill (15 points)}

\noindent\textbf{Exercice 1.} \hfill\textit{(3 points)}\\
Pour relier le château d'eau aux quartiers, le bureau d'études dispose de $9$ tronçons de
canalisation : $5$ tronçons en fonte et $4$ tronçons en PVC, tous distincts. On choisit
\emph{simultanément} $4$ tronçons pour réaliser une conduite.
\begin{enumerate}[label=\arabic*)]
\item Déterminer le nombre de choix possibles de $4$ tronçons parmi les $9$. \hfill\textit{(0,75 pt)}
\item Déterminer le nombre de choix comportant \emph{exactement} $2$ tronçons en fonte et $2$ en PVC. \hfill\textit{(1 pt)}
\item Déterminer le nombre de choix comportant \emph{au moins un} tronçon en PVC. \hfill\textit{(1,25 pt)}
\end{enumerate}

\smallskip
\noindent\textbf{Exercice 2.} \hfill\textit{(3 points)}\\
Trois stations de pompage $S_1$, $S_2$, $S_3$ refoulent de l'eau. En une heure, leurs débits
respectifs $x$, $y$, $z$ (en m\textsuperscript{3}) vérifient le système :
\[ \begin{cases} x+y+z=130\\ 2x+y+3z=290\\ x+2y+z=170 \end{cases} \]
\begin{enumerate}[label=\arabic*)]
\item Résoudre ce système par la méthode du pivot de Gauss. \hfill\textit{(2,25 pt)}
\item En déduire le débit horaire total des trois stations et vérifier le résultat. \hfill\textit{(0,75 pt)}
\end{enumerate}

\smallskip
\noindent\textbf{Exercice 3.} \hfill\textit{(4 points)}\\
La pression résiduelle (en bar) en un point du réseau, $t$ heures après l'ouverture d'une
vanne, est modélisée pour $t\ge0$ par $f(t)=(2-t)\mathrm{e}^{t}+t-2$. On note $(\mathcal{C})$ sa courbe.
\begin{enumerate}[label=\arabic*)]
\item Calculer $f(0)$ et déterminer $\displaystyle\lim_{t\to+\infty}f(t)$. \hfill\textit{(0,75 pt)}
\item Montrer que $f'(t)=(1-t)\mathrm{e}^{t}+1$, puis étudier le signe de $f'$ et dresser le
tableau de variations de $f$ sur $[0,+\infty[$. \hfill\textit{(1,75 pt)}
\item À l'aide d'une intégration par parties, calculer $\displaystyle I=\int_0^{1}(2-t)\mathrm{e}^{t}\dd t$,
puis en déduire $\displaystyle J=\int_0^{1}f(t)\dd t$. \hfill\textit{(1,5 pt)}
\end{enumerate}

\smallskip
\noindent\textbf{Exercice 4.} \hfill\textit{(5 points)}\\
Le plan est muni d'un repère orthonormé $(O\,;\,\vec\imath,\vec\jmath)$. Un bassin de rétention de
la CDE a la forme de l'ellipse $(\mathcal{E})$ d'équation $\dfrac{x^2}{25}+\dfrac{y^2}{9}=1$.
\begin{enumerate}[label=\arabic*)]
\item Déterminer les sommets, les demi-axes et l'excentricité de $(\mathcal{E})$. \hfill\textit{(1,5 pt)}
\item Déterminer les coordonnées des foyers $F$ et $F'$ et les équations des directrices. \hfill\textit{(1,5 pt)}
\item Soit $A(3\,;\,y_A)$ un point de $(\mathcal{E})$ d'ordonnée positive. Calculer $y_A$, puis
montrer par le calcul que $AF+AF'=10$. \hfill\textit{(1,25 pt)}
\item Les vecteurs $\vec u(4\,;\,3)$ et $\vec w(-3\,;\,4)$. Calculer $\vec u\cdot\vec w$ et en
déduire la mesure de l'angle $(\vec u,\vec w)$. \hfill\textit{(0,75 pt)}
\end{enumerate}

\subsection*{Partie B — Évaluation des compétences \hfill (5 points)}

\noindent\textit{La CDE souhaite optimiser l'alimentation en eau du quartier Bota à Limbé.
Elle te sollicite, en tant qu'élève de Terminale~C, pour trois décisions techniques.\\
$\bullet$ Une conduite cylindrique de section circulaire doit transporter un débit fixé : le coût
au mètre est proportionnel au périmètre de la section, et l'on souhaite, à \emph{périmètre}
$P=\SI{2,4}{\metre}$ donné, que l'aire de la section (donc le débit) soit la plus grande possible
en jouant sur le rayon $r$ d'un demi-disque surmontant un rectangle (regard d'entrée).\\
$\bullet$ Le nombre d'abonnés du quartier était de $1500$ en $2024$ et croît de \SI{6}{\percent} par an.\\
$\bullet$ Sur un lot de vannes neuves, \SI{4}{\percent} sont défectueuses ; un technicien en monte $5$ au hasard.}

\begin{enumerate}[label=\textbf{Tâche \arabic*.},leftmargin=2.4em]
\item Le regard a la forme d'un rectangle de base $2r$ et de hauteur $h$ surmonté d'un demi-disque
de rayon $r$. Son périmètre extérieur (les deux côtés verticaux, la base et le demi-cercle) vaut
$P=2h+2r+\pi r=\SI{2,4}{\metre}$. Déterminer le rayon $r$ qui rend l'aire totale maximale. \hfill\textit{(1,5 pt)}
\item Déterminer l'année à partir de laquelle le nombre d'abonnés dépassera $2500$. \hfill\textit{(1,5 pt)}
\item Déterminer la probabilité qu'\emph{au moins une} des $5$ vannes montées soit défectueuse. \hfill\textit{(1,5 pt)}
\end{enumerate}
\par\smallskip{\footnotesize\itshape Présentation générale : 0,5 point.}

% ════════════════════════════════════════════════════
\corrige{du sujet 26}

\noindent\textbf{Partie A — Exercice 1.}
1) Nombre de choix de $4$ tronçons parmi $9$ : $\binom{9}{4}=126$.
2) Exactement $2$ fonte et $2$ PVC : $\binom{5}{2}\binom{4}{2}=10\times6=60$.
3) Au moins un PVC $=$ total $-$ (aucun PVC, c.-à-d. les $4$ en fonte) $=\binom94-\binom54=126-5=\mathbf{121}$.

\noindent\textbf{Exercice 2.}
1) On note $L_1:x+y+z=130$, $L_2:2x+y+3z=290$, $L_3:x+2y+z=170$.
$L_2\leftarrow L_2-2L_1:\ -y+z=30$. $L_3\leftarrow L_3-L_1:\ y=40$.
D'où, dans $-y+z=30$ : $z=30+y=70$. Puis $x=130-y-z=130-40-70=20$.
\textbf{Solution :} $(x\,;\,y\,;\,z)=(20\,;\,40\,;\,70)$.
2) Débit total $=20+40+70=\SI{130}{\cubic\metre}$ par heure (cohérent avec $L_1$).
Vérification dans $L_2$ : $2(20)+40+3(70)=40+40+210=290$ $\checkmark$.

\noindent\textbf{Exercice 3.}
1) $f(0)=(2-0)\mathrm{e}^{0}+0-2=2-2=0$. Quand $t\to+\infty$, $(2-t)\to-\infty$ et $\mathrm{e}^t\to+\infty$,
donc $(2-t)\mathrm{e}^t\to-\infty$ (dominé par l'exponentielle), et $t-2\to+\infty$ ; au total
$f(t)\sim(2-t)\mathrm{e}^t\to-\infty$. Donc $\displaystyle\lim_{t\to+\infty}f(t)=-\infty$.
2) $f'(t)=\big[(2-t)\mathrm{e}^t\big]'+1=\big(-\mathrm{e}^t+(2-t)\mathrm{e}^t\big)+1=(1-t)\mathrm{e}^t+1$.
Sur $[0,+\infty[$ : pour $t=0$, $f'(0)=1+1=2>0$. Posons $g(t)=(1-t)\mathrm{e}^t+1$ ;
$g'(t)=-\mathrm{e}^t+(1-t)\mathrm{e}^t=-t\,\mathrm{e}^t\le0$, donc $g=f'$ est décroissante sur $[0,+\infty[$,
de $f'(0)=2$ vers $\lim_{+\infty}f'=-\infty$. Elle s'annule en un unique $\alpha>0$ avec $f'(\alpha)=0$.
Numériquement $(1-t)\mathrm{e}^t=-1$ donne $\alpha\approx1{,}28$. Donc $f$ croît sur $[0,\alpha]$, décroît
sur $[\alpha,+\infty[$ : maximum en $\alpha$. (Tableau : $f$ part de $f(0)=0$, croît jusqu'à
$f(\alpha)>0$, puis décroît vers $-\infty$.)
3) IPP ($u=2-t$, $u'=-1$, $v'=\mathrm{e}^t$, $v=\mathrm{e}^t$) :
$I=\big[(2-t)\mathrm{e}^t\big]_0^1-\int_0^1(-1)\mathrm{e}^t\dd t=\big[(2-t)\mathrm{e}^t\big]_0^1+\big[\mathrm{e}^t\big]_0^1$.
$\big[(2-t)\mathrm{e}^t\big]_0^1=(1)\mathrm{e}-(2)(1)=\mathrm{e}-2$ ; $\big[\mathrm{e}^t\big]_0^1=\mathrm{e}-1$.
Donc $I=(\mathrm{e}-2)+(\mathrm{e}-1)=2\mathrm{e}-3$.
Puis $J=\int_0^1 f=I+\int_0^1(t-2)\dd t=(2\mathrm{e}-3)+\big[\tfrac{t^2}{2}-2t\big]_0^1
=(2\mathrm{e}-3)+(\tfrac12-2)=2\mathrm{e}-3-\tfrac32=\mathbf{2\mathrm{e}-\tfrac92}\approx0{,}937$.

\noindent\textbf{Exercice 4.}
1) Forme $\frac{x^2}{a^2}+\frac{y^2}{b^2}=1$ avec $a^2=25,\ b^2=9$, donc $a=5$, $b=3$ ($a>b$, grand axe sur $Ox$).
Sommets : $(\pm5\,;\,0)$ et $(0\,;\,\pm3)$. Demi-grand axe $a=5$, demi-petit axe $b=3$.
$c^2=a^2-b^2=25-9=16$, $c=4$ ; excentricité $e=\frac ca=\frac45=0{,}8$.
2) Foyers sur $Ox$ : $F(4\,;\,0)$ et $F'(-4\,;\,0)$. Directrices : $x=\pm\frac ae=\pm\frac{5}{0{,}8}=\pm\frac{25}{4}$.
3) Pour $x=3$ : $\frac{9}{25}+\frac{y_A^2}{9}=1\Rightarrow \frac{y_A^2}{9}=\frac{16}{25}\Rightarrow y_A^2=\frac{144}{25}$,
donc $y_A=\frac{12}{5}$ (ordonnée positive). $A(3\,;\,\frac{12}{5})$.
$AF=\sqrt{(3-4)^2+(\frac{12}{5})^2}=\sqrt{1+\frac{144}{25}}=\sqrt{\frac{169}{25}}=\frac{13}{5}$.
$AF'=\sqrt{(3+4)^2+(\frac{12}{5})^2}=\sqrt{49+\frac{144}{25}}=\sqrt{\frac{1369}{25}}=\frac{37}{5}$.
$AF+AF'=\frac{13}{5}+\frac{37}{5}=\frac{50}{5}=10=2a$ $\checkmark$.
4) $\vec u\cdot\vec w=4\times(-3)+3\times4=-12+12=0$ : les vecteurs sont orthogonaux,
donc l'angle $(\vec u,\vec w)=\frac\pi2$ (90°).

\smallskip
\noindent\textbf{Partie B — Situation.}
\emph{Tâche 1.} Aire totale $=$ rectangle $+$ demi-disque $=2r\,h+\frac12\pi r^2$.
Contrainte $2h+2r+\pi r=2{,}4$, soit $h=\frac{2{,}4-(2+\pi)r}{2}=1{,}2-(1+\tfrac\pi2)r$.
$A(r)=2r\big(1{,}2-(1+\tfrac\pi2)r\big)+\tfrac12\pi r^2=2{,}4r-(2+\pi)r^2+\tfrac\pi2 r^2
=2{,}4r-\big(2+\tfrac\pi2\big)r^2$.
$A'(r)=2{,}4-2\big(2+\tfrac\pi2\big)r=2{,}4-(4+\pi)r$. $A'(r)=0\Rightarrow r=\frac{2{,}4}{4+\pi}\approx\frac{2{,}4}{7{,}1416}\approx\SI{0,336}{\metre}$.
$A''=-(4+\pi)<0$ : c'est bien un maximum. Rayon optimal $r\approx\SI{0,34}{\metre}$
(et $h=1{,}2-(1+\tfrac\pi2)(0{,}336)\approx\SI{0,336}{\metre}$, soit $h=r$).

\emph{Tâche 2.} Nombre d'abonnés l'année $n$ (avec $n=0$ en $2024$) : $u_n=1500\times1{,}06^{\,n}$.
On résout $u_n>2500$, soit $1{,}06^{\,n}>\frac{2500}{1500}=\frac53$, d'où
$n>\frac{\ln(5/3)}{\ln1{,}06}\approx\frac{0{,}5108}{0{,}0583}\approx8{,}77$ : à partir de $n=9$, c'est-à-dire
en \textbf{2033} ($u_9=1500\times1{,}06^9\approx2534$).

\emph{Tâche 3.} Le nombre de vannes défectueuses suit une loi binomiale $\mathcal{B}(5\,;\,0{,}04)$.
$P(\text{au moins une})=1-(0{,}96)^5\approx1-0{,}8154=\mathbf{0{,}185}$, soit environ \SI{18,5}{\percent}.
